% 原文 pp.247–253（印刷 pp.1365–1371）
\section{Lecture 18：四维 $\mathcal{N}=2$ 超对称 Yang--Mills 理论（二）}

\begin{keytakeaways}
本讲的物理图像是：纯 $\mathcal{N}=2$ 超对称 $SU(2)$ Yang--Mills 理论在红外只留下一个 $U(1)$ 矢量多重态，
其真空由手征场 $u=\operatorname{Tr}\phi^2$ 的期望值标记；量子真空模空间 $\mathcal{M}$ 在 $u$ 大处与经典 $u$ 平面重合，
而低能有效耦合 $\tau_{\rm eff}(u)$ 因超对称的全纯性而成为 $\mathcal{M}$ 上的全纯函数。
数学工具主要有三件：$\tau$ 的全纯性（它把单圈以上所有微扰修正都固定下来）、反常的 $U(1)_R$ 对称性
（它把 $\theta$ 角的平移与 $u$ 的相位联系起来，从而定出绕无穷远的单值变换），
以及 Weierstrass 形式的椭圆曲线族 $y^2=x^3+A(u)x+B(u)$ 与全纯二形式 $\eta$ 的周期。
主要结论是：$\mathcal{M}$ 只有一支，而且就是整张 $u$ 平面（可能另加至多两个孤立点）；
$\tau=2i\ln(u/\Lambda^2)/2\pi$ 在微扰论的所有阶精确成立，非微扰修正只能来自瞬子，形如 $\sum_k f_k\Lambda^{4k}/u^{2k}$；
绕无穷远的单值变换是 $SL(2,\mathbb{Z})$ 中的抛物元 $\begin{pmatrix}-1&2\\0&-1\end{pmatrix}$，它使 $\tau\to\tau-2$、$a\to-a$ 并保持 BPS 谱不变；
但整个 $\mathcal{M}$ 上的单值群是非阿贝尔的，故 BPS 谱必在一条实曲线上跳变，而 $\mathcal{M}$ 的奇点正是某个 BPS 超多重态变为无质量之处。
\end{keytakeaways}

% ===== 原文 p.247 =====
\subsection{回顾上一讲的内容}
% 原文 §18.1：Review of material from the last lecture

我们考虑四维 Minkowski 空间中的纯 $\mathcal{N}=2$ 超对称规范理论（即不含超多重态（hypermultiplet）的那种）。规范群取为 $SU(2)$，或与之密切相关的群 $SO(3)$。（事实上，由于我们今天是在可缩的 Minkowski 四维空间上工作，我们将看不出这两个规范群之间的区别。）我们用 $\mathcal{N}=1$ 超对称的记号来写这个理论。构成单个 $\mathcal{N}=2$ 矢量多重态（vector multiplet）的那些场，在四维 $\mathcal{N}=1$ 超 Poincaré 群下分解为一个矢量多重态 $\mathcal{A}$（其中含有规范场）与一个手征超场
\[
\Phi=\phi+\theta\psi+\cdots,
\]
其中 $\phi$ 是（就 Poincaré 群而言的）一个标量场。它是伴随丛的复化（complexification）的一处截面。微观 Lagrange 量为
\[
\mathcal{L}=\int d^4x\,d^2\theta\,\tau_0(\Lambda_0)\operatorname{Tr}W_\alpha W^\alpha+\text{c.c.}+\int d^4x\,d^2\theta\,\operatorname{Im}\tau_0(\Lambda_0)\operatorname{Tr}\bar\Phi\Phi,
\]
其中 $W$ 是超联络（superconnection）$\mathcal{A}$ 的曲率的某一特定部分。这里 $\Lambda_0$ 是为使 Feynman 路径积分有限而引入的截断。我们把 $\tau_0$ 当作 $\Lambda_0$ 的函数来变动，以便重整化这个理论，也就是在 $\Lambda_0\to\infty$ 时使理论中的期望值变为有限。（为了决定如何重整化，我们必须选一个所谓「减除点」（subtraction point）$\Lambda$，在其处固定理论中某一个参数的值。$\Lambda$ 在 $\Lambda_0\to\infty$ 的过程中保持固定，而 $\tau_0$ 则被调节到使理论具有极限。稍后我们将会看到，对这个特定理论存在一个方便的归一化方案，它要牵涉到理论的细节——$u$ 平面等等。但眼下我们先不指明实际所用的重整化方案。）虽然 $\tau_0$ 依赖于截断，它与场无关。

% [Tutor Note: 上式第二项原文印作 $\int d^4x\,d^2\theta\,\operatorname{Im}\tau_0\operatorname{Tr}\bar\Phi\Phi$，这是笔误：$\operatorname{Tr}\bar\Phi\Phi$ 是 Kähler 势型的 D 项，必须对全超空间积分 $\int d^2\theta\,d^2\bar\theta=\int d^4\theta$；只有第一项（规范动能项，被积函数是手征的）才是 $\int d^2\theta$。若照字面用 $d^2\theta$，取到的只是复合算符 $\bar\Phi\Phi$ 的 F 项（形如 $\bar F\phi$ 的项），它既不与第一项构成实的作用量，也给出不了下文 $ds^2$ 里那个 $\operatorname{Im}\tau_0$ 的 Kähler 度量；正确写法是 $\int d^4\theta\,\operatorname{Im}\tau_0\operatorname{Tr}\bar\Phi\Phi$。另须注意约定：由 (18.1) 的 $\operatorname{Im}\tau=4\pi/g_{\rm eff}^2$ 与 p.251 的 $\theta=2\pi\operatorname{Re}\tau$ 可知，本文取 $\tau=\theta/2\pi+4\pi i/g^2$，而非另一种常见的 $\tau=\theta/2\pi+i/g^2$；若换用后者，(18.1)、(18.2)、(18.5) 以及瞬子因子 $\exp(2\pi ik\tau)$ 中所有 $4\pi$ 都要改写，这是核对本讲公式时最容易出错的地方。]

我们今天的目的是求解这个理论：显式地理解量子真空的模空间（moduli space）$\mathcal{M}$，并得出理论在每个真空态附近的描述。正如我们上次所确立的，这个理论的经典真空空间

% ===== 原文 p.248 =====
由 $\mathfrak{su}(2)$ 中常数矩阵的规范等价类构成，即
\[
\frac{1}{\sqrt2}\begin{pmatrix}a&0\\0&-a\end{pmatrix},
\]
模掉 Weyl 群的作用；Weyl 群是二阶群，它把 $a$ 与 $-a$ 互换。这个空间解析同构于 $\mathbb{C}$，而在这张平面上应当使用的自然局部全纯坐标是 $u=\operatorname{Tr}\phi^2=a^2$。经典真空的解析空间虽然是光滑的，但在经典真空的模空间中有一个点使 Kähler 度量奇异，这就是原点。奇异性是这样一件事实的后果：在这个点上存在着更多的无质量场（事实上低能有效理论是一个超对称的 $SU(2)$ 理论）。此外，正如我们上次讨论过的，在离开 $u=0$ 处，低能有效理论是一个纯 $\mathcal{N}=1$ 超对称 $U(1)$ 规范理论。于是，在离开 $u=0$ 处，这个有效理论的性质由带 puncture（打洞）的 $u$ 平面上的一族环面来描述。就经典真空的模空间而言，这是一族具有常数 $j$ 不变量的环面，其 $j$ 不变量由 Lagrange 量中 $\tau_0$ 的选取给出。对无穷远点，这一族环面的单值变换为 $-1$。$u$ 平面上的 Kähler 度量为
\[
ds^2=\operatorname{Im}\tau_0\,\frac{du\,d\bar u}{4|u|^2}.
\]
（注意：度量在 $u=0$ 处有奇异性，这与经典上该真空处存在额外无质量场这一事实相符。）

% [Tutor Note: 「离开 $u=0$ 处低能有效理论是纯 $\mathcal{N}=1$ 超对称 $U(1)$ 规范理论」中的 $\mathcal{N}=1$ 应为 $\mathcal{N}=2$ 之误。判据有三：(i) 本讲 p.247 开篇即声明研究对象是纯 $\mathcal{N}=2$ 理论；(ii) 第 17 讲（p.245）在讨论同一模型时明言「在原点有效理论不是纯 $\mathcal{N}=2$ 超对称 $U(1)$ 规范理论」，反过来说离开原点就是；(iii) 紧接着的句子说该有效理论由常数 $j$（即 $\tau_0$）的一族环面描述，而 $\tau_0$ 正是 $\mathcal{N}=2$ $U(1)$ 矢量多重态的复耦合。译文照原样保留，读者应自行改读。]

\subsection{关于量子真空模空间 $\mathcal{M}$ 的首批结果}

设 $\mathcal{M}$ 是我们的 $\mathcal{N}=2$ 超对称纯规范理论的量子真空空间。利用微扰论与渐近自由，在 $\mathcal{M}$ 上无穷远附近（即 $u$ 大处），低能谱恰好由一个 $U(1)$ 矢量多重态构成；我们将假定这就是 $\mathcal{M}$ 上一般点处正确的低能谱。若果真如此，就可以搬用上周的论证。量子真空的空间具有一种相当特殊的 Kähler 度量，并且在这个真空空间之上有一个全纯的环面族。

场 $u$ 是经典真空空间上的一个坐标，它是理论中的一个手征场，因而它在各个量子真空中的期望值定义了 $\mathcal{M}$ 上的一个全纯函数。我们把这个函数记为 $u:\mathcal{M}\to\mathbb{C}$：
\[
u(\Omega)=\langle\Omega|u|\Omega\rangle,\qquad\text{对一切 }\Omega\in\mathcal{M}.
\]

$\mathcal{M}$ 中与经典真空的 $u$ 平面关系最密切的部分，是 $u$ 很大的那个区域。更精确地说，令 $\Lambda$ 为理论的重整化标度，并在下面的讨论中令它固定。假定我们研究的是某个经典真空 $\Omega$ 附近、且 $u(\Omega)\gg\Lambda^2$ 的理论。（注意：由于 $u$ 的量纲是质量的平方，它自然地与 $\Lambda^2$ 相比较。）由于我们的 $SU(2)$ 理论是渐近自由的，而它所趋近的 $U(1)$ 理论在红外是自由的，所以若 $u\gg\Lambda^2$，量子理论在 $\Omega$ 附近就被经典理论很好地近似。这意味着 $\Omega$ 附近的低能量子理论由 $u$ 平面上的一个 $\sigma$ 模型描述，或许还带有某个超势的微扰。但 $\mathcal{N}=2$ 超对称意味着任何这样的超势都只能是常数，因此不能在 $\Omega$ 附近抬高真空的简并度。（此外还有 Fayet--Iliopoulos $\theta$ 项的可能性，但 $SU(2)_R$ 对称性把这些都排除了。）这个论证所表明的是：对大的 $u$，每一个经典真空对应唯一的量子真空。

% ===== 原文 p.249 =====
也就是说，存在常数 $C(\Lambda)>0$ 与开子集 $\mathcal{U}\subset\mathcal{M}$，使得 $u:\mathcal{U}\to\{z\in\mathbb{C}\mid |z|>C\}$ 是全纯同构。

在第 16 讲中，以及更详细地在 [II-N=1 Dynamics] 中，我们已经见过另一些规范理论的例子：在经典真空空间的无穷远附近，规范群 $SU(N_c)$ 破缺到一个非阿贝尔规范群 $SU(N_c-N_f)$。由于这个较小群对应的规范理论在红外不是自由的，我们不能断言量子真空空间在无穷远附近可以等同于经典真空空间上某个超势的临界点集合。事实上，在这些例子里，量子真空的模空间是经典真空空间在无穷远附近的一个有限覆盖空间。再强调一次，这是可能的，因为低能理论在红外不是自由的。而当低能理论是自由的时候——正如现在所讨论的情形——并且当 $\mathcal{N}=2$ 超对称排除了超势的存在时，经典模空间与量子模空间可以在经典真空空间的无穷远邻域上等同起来。

我们将不作任何论证地假定：模空间 $\mathcal{M}$ 在适当意义下是完备的（complete），从而它对应于一个开的复代数簇。

通过考虑经典极限，我们找到了 $\mathcal{M}$ 的一个连通分支，它有一个端（end），在 $u$ 之下同构于 $\mathbb{C}$ 中无穷远的一个邻域。于是有两个问题摆在我们面前：
\begin{itemize}
\item $\mathcal{M}$ 是否还有别的「端」，即经典上看不到的端？
\item $\mathcal{M}$ 是否还有别的不可约分支，是经典上看不到的？
\end{itemize}

对第一个问题的回答是「没有」。量子模空间的所有端在经典层面都是可见的，也就是说，它们与经典模空间的端相联系。这是关于量子场论的一个一般性论断：量子效应是有界的，因而不能在从经典真空空间过渡到量子真空空间的过程中造成非逆紧（non-proper）的改变。

至于第二个问题，还没有已知的超对称理论例子，其真空空间的量子版本存在经典上看不到的分支。不过，在这样一般的层次上，尚未证明这样的分支不存在。在我们的特殊情形中，我们可以给出论证表明没有别的分支。首先要指出的是，由对第一个问题的回答，$\mathcal{M}$ 的任何别的分支都必须是紧的。为了排除其他紧分支的可能性，我们利用这个理论存在一个相关（relevant）的 $\mathcal{N}=1$ 超对称微扰这一事实。这个微扰就是
\[
\Delta\mathcal{L}=\int d^4x\,d^2\theta\,\varepsilon u+\text{c.c.},
\]
其中 $\varepsilon$ 是复常数。正如我们在第 16 讲 §16.2 中所见（也见 [II-N=1 Dynamics, §2.2]），加上这个微扰之后，所得的 $SU(2)$ 理论有两个真空，人们相信其中每一个都有质量间隙与禁闭。这个结果必须以某种方式在低能有效描述中得到重现。可是，沿着一个假设中的、经典上看不到的紧分支，$\Delta\mathcal{L}$ 在低能有效描述中会是什么样子呢？由 $\mathcal{N}=1$ 超对称，它必须看起来像一个有效超势（外加带导数、不能抬高真空简并度的项）；但是，作为紧分支上的全纯函数，这个超势将是常数，因而它是平凡的。所以即便把 $\Delta\mathcal{L}$ 加进 Lagrange 量，紧分支上的真空也会存留下来，于是

% ===== 原文 p.250 =====
低能理论就会有过多的真空态——只要那个紧分支具有正维数。（这个论证仍然允许紧分支由孤立点组成——至多两个这样的点——但我们稍后将通过数奇点的方法看到，这种情况同样不会发生。）

至此我们已经论证：$\mathcal{M}$ 只有一个分支（可能再加上至多两个孤立点），并且这个分支只有一个无穷远邻域，它被 $u$ 同构地映到 $u$ 平面中无穷远的一个邻域。按我们的完备性假定，$\mathcal{M}$ 是一条可能奇异的、不可约的复曲线，在 $u=\infty$ 处有一个 puncture。只具有一个极点、且该极点是 $\mathcal{M}$ 的无穷远处的单极点的全纯函数 $u$ 的存在，蕴涵着 $\mathcal{M}$ 实际上就是复平面，而 $u$ 是 $\mathcal{M}$ 上的整体全纯参数。

% [Tutor Note: 这一步作者写得很简，省掉的是「度数一的亚纯函数给出到 $\mathbb{P}^1$ 的同构」这一基本事实。补足如下：把带 puncture 的 $\mathcal{M}$ 添上 $u=\infty$ 一点紧化为紧 Riemann 面 $\bar{\mathcal{M}}$，则 $u$ 是 $\bar{\mathcal{M}}$ 上的亚纯函数，且只有一个单极点，故其次数（极点数，计重数）为 $1$；而度数 $1$ 的亚纯函数给出双全纯映射 $\bar{\mathcal{M}}\cong\mathbb{P}^1$。再去掉 $u=\infty$ 这一点即得 $\mathcal{M}\cong\mathbb{C}$。注意「单极点」这一条不可省：若允许二阶极点，$\mathcal{M}$ 就可以是 $\mathbb{P}^1$ 的二重覆盖，即 $u$ 是从一条超椭圆曲线上的投影，结论立刻不成立。]

沿用上一讲的分析，我们的目标是通过求出 $\mathcal{M}$ 上的椭圆曲线族（取 Weierstrass 形式）
\[
y^2=x^3+A(u)x+B(u)
\]
以及这个族全空间上的一个全纯二形式 $\eta$
\[
\eta=f(u)\,du\,\frac{dx}{y}
\]
来求解这个理论。椭圆曲线的复结构将确定 $\tau(u)$，而 $\eta$ 的周期则确定对偶的局部全纯坐标的微分 $da$ 与 $da_D$。在 $\mathcal{M}$ 的奇点处，椭圆曲线可能变得奇异，$\tau(u)$ 可能发散。

\subsection{$\mathcal{M}$ 中无穷远处的性质}

正如我们已经提到的，利用 $SU(2)$ 理论的渐近自由，我们可以得到 $\mathcal{M}$ 中无穷远附近（即 $u$ 大处）的信息。这里的目标是计算量子真空模空间上椭圆曲线族在无穷远处的单值变换。我们先研究经典真空空间的无穷远处的单值变换。回忆一下，经典真空空间就是矩阵
\[
\phi=\frac{1}{\sqrt2}\begin{pmatrix}a&0\\0&-a\end{pmatrix}
\]
的空间，模去 $a\mapsto-a$。对偶坐标 $a_D$ 由 $a_D=\tau_{\rm cl}a$ 给出，其中 $\tau_{\rm cl}$ 是常数（前面记作 $\tau_0$），具有很大的虚部。$u$ 平面上的整体坐标是 $u=\operatorname{Tr}\phi^2=a^2$，于是 $a=\sqrt u$。由此可以清楚地看出，在这个经典族中，绕无穷远的单值变换是 $-1$。

对于量子真空的空间 $\mathcal{M}$，绕无穷远的单值变换是不同的。我们给出两个计算它的论证：第一个基于渐近自由，第二个基于 $U(1)_R$ 反常。为了计算 $u$ 大时的 $\tau_{\rm eff}(u)$，我们考虑单圈公式。适当地选取质量参数 $\Lambda$ 之后，公式为
\begin{equation}
\tag{18.1}
\label{eq:18-1}
\operatorname{Im}\tau=\frac{4\pi}{g_{\rm eff}^2}=4\pi\Bigl(2b_0\ln\bigl(\sqrt{|u|}/\Lambda\bigr)\Bigr),
\end{equation}
其中 $b_0$ 是单圈 $\beta$ 函数系数。

单圈计算给出
\[
b_0=\frac{1}{4\pi^2}.
\]

% ===== 原文 p.251 =====
由于 $\tau$ 是全纯的，知道它的虚部就唯一地确定了它。我们看到，对大的 $u$ 有
\begin{equation}
\tag{18.2}
\label{eq:18-2}
\tau_{\rm eff}=\frac{2i\ln(u/\Lambda^2)}{2\pi}.
\end{equation}
于是我们看到，绕无穷远的单值变换把 $\tau$ 平移 $-2$，从而把 $\theta$ 平移 $-4\pi$。这意味着 $\mathcal{M}$ 中绕无穷远的单值变换由
\[
\begin{pmatrix}-1&2\\0&-1\end{pmatrix}
\]
给出。

% [Tutor Note: 由 (18.1) 到 (18.2) 的算术值得逐项核对，这是本讲最容易「差一个 $2$」的地方。取 $a=\sqrt u$，则 $\ln(\sqrt{|u|}/\Lambda)=\ln(|a|/\Lambda)$，而 $\ln(|u|/\Lambda^2)=2\ln(|a|/\Lambda)$。于是 (18.1) 给出 $\operatorname{Im}\tau=4\pi\cdot 2b_0\ln(|a|/\Lambda)$，代入 $b_0=1/4\pi^2$ 得 $\operatorname{Im}\tau=\frac{2}{\pi}\ln(|a|/\Lambda)$；(18.2) 给出 $\operatorname{Im}\tau=\frac{1}{\pi}\ln(|u|/\Lambda^2)=\frac{2}{\pi}\ln(|a|/\Lambda)$，二者完全一致。这个一致性同时验证了 $\beta$ 函数的归一化：把 $\alpha=1/g^2$ 的跑动写成 $d\alpha/d\ln\mu=2b_0$（即 $\beta(g)=-b_0g^3$），正是 $b_0=1/4\pi^2$。还要注意 (18.2) 中被取对数的是 $u$ 而不是 $|u|$：$\tau$ 必须对 $u$ 全纯（这正是它能被椭圆曲线的复结构决定的前提），所以 $|u|$ 只能出现在取虚部之后。代价是 $\ln u$ 有分支：$u\to e^{2\pi i}u$ 时 $\ln u\to\ln u+2\pi i$，$\tau\to\tau+\frac{2i\cdot 2\pi i}{2\pi}=\tau-2$，这正是下文单值变换的来源。]

换一种说法，我们考虑的是如下系统：当我们绕着 $u$ 平面中的无穷远转一圈时，
\begin{equation}
\tag{18.3}
\label{eq:18-3}
\begin{aligned}
u^2&=a,\\
\frac{da_D}{da}&=\tau=\frac{2i\ln(u/\Lambda^2)}{2\pi}.
\end{aligned}
\end{equation}
显然，$a$ 回到 $-a$，而 $\tau$ 回到 $\tau-2$。上面那个单值矩阵的形式于是很容易得出。

% [Tutor Note: (18.3) 第一行原文印作 $u^2=a$，这是笔误，应为 $u=a^2$（等价地 $a=\sqrt u$）。依据：本页上文与 p.250 都写 $\phi=\frac{1}{\sqrt2}\operatorname{diag}(a,-a)$、$u=\operatorname{Tr}\phi^2=a^2$；而且紧接着的句子「显然 $a$ 回到 $-a$」只对 $a=\sqrt u$ 成立——若真按 $u^2=a$，则 $u\to e^{2\pi i}u$ 时 $a\to e^{4\pi i}a=a$，根本不变号，$\frac{da_D}{da}$ 也就无从得到 $\tau-2$。译文照原样保留，读者应自行改读。（同理，第二行右端 $2i\ln(u/\Lambda^2)/2\pi$ 与 (18.2) 一致，可用来校验第一行。）]

% [Tutor Note: 单值矩阵写成 $\begin{pmatrix}-1&2\\0&-1\end{pmatrix}$ 时，必须明确它作用在列向量 $(a_D,a)^T$ 上，即 $a_D\to-a_D+2a$、$a\to-a$；此时 $\tau=da_D/da\to(-da_D+2da)/(-da)=\tau-2$，与「$\tau$ 平移 $-2$」相符。若改用 $(a,a_D)^T$ 的顺序，矩阵就变成 $\begin{pmatrix}-1&0\\2&-1\end{pmatrix}$，$2$ 挪到左下角。抄写时把 $2$ 放在右上角而心中想的却是 $(a,a_D)$ 的顺序，会得到 $\tau\to\tau+2$ 的错误结论。另注意此矩阵是 $SL(2,\mathbb{Z})$ 中的抛物元 $-(T^{-2})$，它生成 $\tau\to\tau-2k$ 的平移，本身并不返回原状（其平方为 $\begin{pmatrix}1&-4\\0&1\end{pmatrix}$）；这与经典情形中真正等于 $-1$ 的单值变换形成对照。]

现在让我们用反常的 $U(1)_R$ 对称性重新计算这个单值变换，看看一切是怎样自洽起来的。经典地，规范多重态中各分量场的 $U(1)_R$ 荷如下：$A$ 的荷为零，$\lambda,\psi$ 的荷为 $1$，$a$ 的荷为二。这意味着 $u=a^2$ 的荷为 $4$。（这与 Donaldson 理论中 $\mu(\text{pt})$ 是 $4$ 维的这一事实有关。）这意味着，在把手征超空间坐标乘以 $e^{i\beta}$ 的那个 $U(1)_R$ 对称性下，变量 $u$ 按 $u\mapsto ue^{\pm4i\beta}$ 变换。事实上，符号是负的。另一方面，如我们所知，存在一个量子反常，把这个 $U(1)_R$ 破缺为 $\mathbb{Z}/8\mathbb{Z}$。这是一个指标计算，用到了这样一件事实：理论中有 $2$ 个费米场（按上周的记号即 $\lambda$ 与 $\psi$）位于伴随表示中，而 $\mathfrak{su}(2)$ 的伴随表示的 Casimir 是定义表示 Casimir 的 $4$ 倍。$2\times4=8$，所以瞬子场中的反常总是可被 $8$ 整除，对称性便是 $\mathbb{Z}/8\mathbb{Z}$。（本质相同的计算也确定了三流形的 Floer 群是 $\mathbb{Z}/8\mathbb{Z}$ 分次的这一事实。）这意味着，在 $U(1)_R$ 对称性 $e^{i\beta}$ 之下，$\theta$ 角 $\theta=2\pi\operatorname{Re}\tau$ 变换为 $\theta-8\beta$。于是，当我们在 $u$ 平面中绕无穷远走一圈时，$u$ 的辐角改变 $2\pi$，相应地有一个 $2\pi/4=\pi/2$ 的 $U(1)_R$ 变换，而 $\theta$ 减少 $4\pi$。这意味着在绕无穷远的单值变换下 $\tau\mapsto\tau-2$。

当然，这些并不是彼此独立的计算。超对称性在 $U(1)_R$ 反常与 $\beta$ 函数之间强加了一个关系，从而保证一切都自洽。

事实上，关于无穷远处 $\tau$ 的性质我们还可以多说一些。对 $|u|\gg\Lambda^2$，在微扰论中我们有级数
\begin{equation}
\tag{18.4}
\label{eq:18-4}
\tau(u)=\frac{2i\ln(u/\Lambda^2)}{2\pi}+\tau_{\ge2},
\end{equation}
其中第一项是经典值与（上面算出的）单圈修正之和，第二项是更高圈修正之和。这些更高圈的修正给出 $g^2$ 的幂级数（$k$ 圈修正给出 $g^{2k-2}$ 的幂）。由于 $\theta=\operatorname{Re}\tau$ 没有圈修正，又由于 $\tau$ 是 $u$ 的全纯函数，立刻可以推出所有更高圈的修正都为零：$\tau_{\ge2}=0$。因此，作为微扰论中的陈述，(18.2) 式在所有阶上都成立。然而，对微扰

% ===== 原文 p.252 =====
展开还存在非零的修正，它们来自瞬子。当然，这些在微扰论中是看不到的。让我们研究它们的性质（使用我们在第 16 讲见过的某类论证）。$\tau$ 的全纯性蕴涵着：对微扰论的修正应当来自在一半超对称变换下不变的场，也就是来自瞬子。一个 $k$ 瞬子场的作用量为 $8\pi^2k/g^2$，所以 $k$ 瞬子振幅带有一个因子 $\exp(-8\pi^2k/g^2)e^{ik\theta}=\exp(2\pi ik\tau)$。重整化群蕴涵着这个因子必须伴随一个 $\Lambda^{4k}$ 的因子。由量纲分析，这就必须是 $\Lambda^{4k}/u^{2k}$。作为检验，我们注意到（正如第 16 讲中讨论过的类似问题那样），$k$ 瞬子场以 $8k$ 个单位破坏 $U(1)_R$ 对称性；由于 $u$ 的荷为 $4$，这就解释了 $u^{-2k}$ 这个因子。于是，$k$ 瞬子对 $\tau$ 的贡献具有形式
\[
f_k\exp(-8\pi^2k/g^2)e^{ik\theta}=f_k\Lambda^{4k}e^{2\pi ik\tau}=f_k\frac{\Lambda^{4k}}{u^{2k}},
\]
其中 $f_k$ 是无量纲常数；而 $\tau(u)$ 对 $|u|\gg1$ 的完整展开具有形式
\begin{equation}
\tag{18.5}
\label{eq:18-5}
\tau(u)=\frac{2i\ln(u/\Lambda^2)}{2\pi}+\sum_{k\ge1}f_k\frac{\Lambda^{4k}}{u^{2k}}.
\end{equation}

% [Tutor Note: 中间那个等号 $\exp(-8\pi^2k/g^2)e^{ik\theta}=\exp(2\pi ik\tau)$ 只有在 $\tau=\theta/2\pi+4\pi i/g^2$ 下才成立：$2\pi ik\tau=ik\theta+2\pi ik\cdot\frac{4\pi i}{g^2}=ik\theta-\frac{8\pi^2k}{g^2}$。若用另一套常见约定 $\tau=\theta/2\pi+i/g^2$，指数上就会缺因子 $4\pi$，瞬子因子与单圈公式 (18.1) 也不再匹配。另外，末句「对 $|u|\gg1$ 的完整展开」在量纲上应是 $|u|\gg\Lambda^2$（与 p.251 的写法一致）：级数的展开参数是 $\Lambda^{4k}/u^{2k}$，要求它小而不仅是 $u$ 数值大。若把 $\Lambda$ 取作质量单位、$u$ 以 $\Lambda^2$ 为单位，两种写法才等价。]

\subsection{BPS 态与 $\mathcal{M}$ 中的奇点}

其次，我们想理解 $\mathcal{M}$ 中的奇点。所谓奇点，我们指模空间中物理发生改变的点。这很可能是 Kähler 度量或其它低能耦合出现某种奇异性的点。作为一个普遍原则，我们预期奇点必定与粒子变为无质量有关。在此情形下，由于这些粒子会给低能极限（它是一个 $U(1)$ 理论）添上奇点，它们应当在 $U(1)$ 下带荷，这是合理的。变为无质量的粒子当然会组成 $\mathcal{N}=2$ 超多重态。更具体地说，若变为无质量的粒子自旋至多为 $1/2$，它们就必须组成一个 BPS 饱和的超多重态，因为这是所有自旋 $\le1/2$ 的 $\mathcal{N}=2$ 多重态中唯一的类型。下面我们给出理由，使人相信理论不会获得额外的、自旋为 $1$ 的无质量粒子；而正如我们在第 2 讲中学到的，它不可能有自旋大于 $1$ 的无质量粒子。所以我们将试图用 BPS 饱和超多重态变为无质量来解释这些奇点。BPS 饱和的超多重态，若并非永远无质量，就必须带有非零的中心荷；在目前这个理论中（正如我们上周讨论的），中心荷就是电荷与磁荷，因此所关心的 BPS 饱和超多重态必须是带电或带磁荷的。

在带 $\mathcal{N}=2$ 超对称的纯 $SU(2)$ 规范理论中，$u$ 平面无穷远附近有两类 BPS 态。我们有带电的矢量多重态，即与 $SU(2)$ 自发破缺部分相联系的规范玻色子。它们有 $n_e=\pm1$ 与 $n_m=0$。\footnote{$n_e$ 与 $n_m$ 当然是 $U(1)$ 生成元 $N_e$ 与 $N_m$ 的本征值。} 我们还有磁单极子，它们有 $n_m=\pm1$ 而 $n_e$ 取任意值。注意这个 BPS 谱在无穷远处的单值变换下不变，那里的变换是 $n_e\mapsto-n_e+2n_m$ 与 $n_m\mapsto-n_m$。磁单极子组成 BPS 饱和的超多重态。（「dyon」这个词有时用于指同时带电荷与磁荷的粒子，所以这些粒子大多是 dyons。）

% ===== 原文 p.253 =====
我们断言：这个谱虽然在无穷远处的单值变换下不变，却不可能在我们这个族的整个单值群下不变。理由是：$\mathcal{M}$ 整个上的单值变换的像，是 $SL(2,\mathbb{Z})$ 的一个非阿贝尔子群（否则 $\operatorname{Im}\tau$ 不可能处处为正）；而上面描述的那个 BPS 谱在这样的非阿贝尔群下并不不变。这意味着 BPS 谱不可能在 $\mathcal{M}$ 上处处连续，也就是说必定存在一条跳变轨迹（jumping locus），即 $\mathcal{M}$ 中的一条实曲线，沿着它额外的 BPS 态变为无质量。正如我们在上一讲所见，这发生在 $a/a_D$ 变为实的时候。

% [Tutor Note: 「否则 $\operatorname{Im}\tau$ 不可能处处为正」这句括号里的理由，作者写得极简，实为模曲线的基本事实，且只是提示性的论证，值得补足：(i) $\tau$ 取值于上半平面 $\mathbb{H}$，单值群 $\Gamma\subset SL(2,\mathbb{Z})$ 按 $\tau\mapsto\frac{a\tau+b}{c\tau+d}$ 作用；(ii) $SL(2,\mathbb{Z})$ 的无限阿贝尔子群（除中心 $\pm I$ 外）必由抛物元生成、因而固定某个尖点，$\mathbb{H}/\Gamma$ 就退化为一个尖点的邻域，$\tau$ 在该邻域上单值，这正意味着「谱无需跳变」——与存在多个不同类型的退化（不同奇点）这一事实矛盾。故 $\Gamma$ 必须非阿贝尔。(iii) 至于「非阿贝尔 $\Gamma$ 不保持该谱」：$\Gamma$ 在荷格 $\mathbb{Z}^2$ 上的轨道对无限非阿贝尔子群通常是无限的，而每一固定点处 BPS 谱所对应的荷是离散集，二者不能一致；于是沿闭合路径回到同一点时谱必须改变，即必在某条实曲线上跳变。这里「$\Gamma$ 非阿贝尔」是必要条件而非充分条件的严格证明，可参 Seiberg--Witten 原文 (Nucl. Phys. B 426 (1994) 19) 及关于 marginal stability 的后续文献。另，判据「$a/a_D$ 实」来自中心荷 $Z=n_e a+n_m a_D$ 的相位：两个态同时变得无质量（$Z$ 同相位）的条件正是 $a/a_D\in\mathbb{R}$。]

我们已经把 $\mathcal{M}$ 等同于复 $u$ 平面。接下来我们想分析 $\mathcal{M}$ 的奇点。$\mathcal{M}$ 中必定有奇点，因为无穷远处的单值变换是非平凡的。如我们已经论证的，$\mathcal{M}$ 中的奇点与某个超对称多重态相联系，该多重态由带电的 BPS 态变为无质量而构成。我们有两类 BPS 多重态——多重态中最大自旋为 $1$ 的那些（矢量多重态），以及最大自旋为 $1/2$ 的那些（超多重态）。一个矢量多重态变为无质量，对应于一个非阿贝尔规范群（$SU(2)$）的恢复。有几个理由使人不期待这件事发生。首先，在 $SU(2)$ 未破缺时，理论在红外是不稳定的。经典模空间的奇点处发生的正是这件事。由于 $\beta$ 函数为负，我们已经从那个理论流走了，所以不该指望它会出现在我们假定在红外为自由的那个低能极限中。当然，如果我们的规范理论有足够多的超多重态变为无质量，使 $\beta>0$，那么我们也许会看到这样一个低能极限。但在那些点上，变为无质量的多重态的最大自旋同时有 $1$ 和 $1/2$。所以我们能做的最简单的事情，就是假定只有一个最大自旋为 $1/2$ 的超多重态变为无质量。\footnote{在二维情形，人们有 Zamolodchikov 的 $c$ 定理——第 13 讲中提到过——它给出了一个确切的意义，说明红外比紫外有更少的无质量自由度。人们猜测类似的定理在任何维数都成立，但尚未证明。如果在四维有这样一个定理，人们就可能对奇点处可以有多少、什么样的无质量场得到一个确切的事先（a priori）界。这样的界很可能会排除关于这个理论可能行为的多数离奇假设。} 让我们分析这种可能性。

假定 $\mathcal{M}$ 的奇点与一个最大自旋为 $1/2$ 的多重态变为无质量相联系。回忆一下：一个最大自旋为 $1/2$ 的有质量 $\mathcal{N}=2$ 超对称多重态自动是 BPS 饱和的。因此我们也许期待，正在发生的事情是：理论中在 $|u|\gg\Lambda^2$ 时存在的那些 BPS 态之一，在奇点处变为无质量。当然，也可能发生更复杂的现象。例如，当我们在 $\mathcal{M}$ 中从无穷远向内走时，可能在到达奇点之前就穿过 BPS 谱的一条跳变线，于是在奇点处变为无质量的可能是一个「新的」BPS 饱和态，而不是 $u=\infty$ 附近原有的那些态之一。但我们最终会看到，产生这个奇点的态正是无穷远处的 BPS 态之一。

让我们假定某个超多重态变为无质量。把它的荷记为 $(n_m,n_e)$。我们已经论证过这个超多重态必须是带荷的，使得 $q=\mathrm{g.c.d.}(n_m,n_e)$ 是一个定义良好的正整数。在把理论用一种 Lagrange 量表示出来、并相差一个 $SL(2,\mathbb{Z})$ 变换之后，我们可以安排 $n_e=q$ 而 $n_m=0$。于是，在奇点 $u_0\in\mathcal{M}$ 附近，我们的理论是一个 $U(1)$ 理论，带有一个在 $u=u_0$ 处变为无质量、荷为 $q$ 的粒子。（但若 $n_m\neq0$，那么在 $u=u_0$ 附近理论表示中的 $U(1)$ 规范场就与

% （本页正文在此处中断，该句续于 p.254；由于本片段截止于 p.253，此处按页末原样收束。）
